\subsection{Integration}
\subsubsection{Stammfunktionen, unbestimmtes Integral}

Sei $I \subset \mathbb{R}$ ein Intervall und $f(x):I \rightarrow \mathbb{R}$ eine Funktion. Eine differenzierbare Funktion $F:I \rightarrow \mathbb{R}$ mit der Eigenschaft
$$F'(x)=f(x), x \in I$$
nennt man die \definition{Stammfunktion} von $f$ auf dem Intervall $I$.

Das Bestimmen von $F(x)$ aus $f(x)$ nennt man \definition{(un)bestimmte Integration.}

\textit{\textbf{Beobachtung:} Stammfunktion \greenBox{ist nicht eindeutig}.}

Ist $F(x)$ Stammfunktion, so ist auch \greenBox{$F(x)+C$} für eine beliebige Konstante $C$ eine Stammfunktion.

$f(x)$ hat eine ganze Familie von Stammfunktionen: \definition{unbestimmtes Integral $\int f(x) \, dx$}. (hier ist $f$ \definition{Integrand}).

\subsubsection{Rechenregeln}

\includegraphics[width=0.75\columnwidth]{04_funktionen/integral_add_mult}

\includegraphics[width=0.9\columnwidth]{04_funktionen/integral_weitere}

$$\int \ln(x)\,dx = x \cdot \ln(x) - x + C$$

\theorem{Partielle Integration}

\textbf{Herleitung:} 
$f(x) \cdot g(x) = \int \frac{d}{dx}(f(x) \cdot g(x))\,dx=\int f'(x) \cdot g(x) + f(x) \cdot g'(x) \, dx = \int f'(x) \cdot g(x) \, dx + \int f(x) \cdot g'(x) \, dx$ 
\textit{(Integration ist Umkehrfunktion von Ableitung; Produkteregel; danach umstellen)}

$$f'(x) \cdot g(x) \, dx = f(x) \cdot g(x) - \int f(x) \cdot g'(x) \, dx$$

$f'(x)$ ist Funktion, die gut integriert werden kann; $g(x)$ ist Funktion, die gut abgeleitet werden kann.

\textit{L (Logarithmus), I (inverse Winkelfunktionen $\arcsin, ...$), A (algebraische Funktionen $x^2, 5x^3, ...$), T (Trigonometrische Funktionen $\sin, \dots$), E (Exponentialfunktionen $e^x, 5a^x, \dots$); Fkt. am Anfang vereinfachen sich durch Ableiten $\implies$ als $g(x)$ wählen.}

\theorem{Integration durch Substitution}

\textbf{Herleitung:}
$f(g(x))+C = \int \frac{d}{dx} (f(g(x)))\,dx = \int f'(g(x)) \cdot g'(x)\,dx = \int \frac{df}{dy} \frac{dy}{dx}\,dx = \int \frac{df}{dy}\, dy, y=g(x)$
\textit{(Integration ist Umkehrfunktion von Ableitung; Kettenregel; $y=g(x)$ ersetzen)}

$$\int f'(g(x))\cdot g'(x)\,dx = \int \frac{df}{dy}\, dy = f(g(x))+C$$

\textbf{Trigonometrische Substitution}
Für Integrale der Form $\int (a \pm x^2)^\alpha \, dx$, $\alpha = \pm 1, \pm \frac{1}{2}$

$\sqrt{1-\sin^2}=\sqrt{\cos^2} = \cos$
\begin{itemize}
    \item $\sin^2(x)+\cos^2(x)=1$
    \item $1+\tan^2(x)=\frac{1}{\cos^2(x)}$
    \item $\cosh^2 - \sinh^2=1$
\end{itemize}



Formen $\int \sqrt{Ax^2 + B x + C} \, dx$ mit quadratischem Ergänzen.

\subsection{Riemann-Integral}
\textbf{Idee:} Unterteile Funktion in $n$ (nicht zwingend gleich grosse) Teilintervalle, wähle jeweils einen repräsentativen Funktionswert in diesem Intervall, summiere alle einzelnen "Teilflächen".

Danach $n \rightarrow \infty$, sodass die Länge der Teilintervalle gegen $0$ geht.

Sei $[a,b] \subset \mathbb{R}$ ein kompaktes Intervall.

\definition{Zerlegung} von $[a,b]$ ist endliche Menge von Punkten $a=x_0 < x_1 < x_2 < ... < x_{n-1} < x_n = b$ mit $n \in \mathbb{N}$.

$x_i$ heissen \definition{Unterteilungspunkte.}

Zerlegung führt zu \definition{Partition} des Intervalls $[a,b] = \{x_0\} \cup (x_0, x_1) \cup \{x_1\} \cup (x_1, x_2) \cup \dots \cup (x_{n-1}, x_n) \cup \{x_n\}$.

Eine Zerlegung $a=y_0 < y_1 < \dots < y_{m-1} < y_m = b$ ist \definition{Verfeinerung} der Zerlegung $a=x_0 < x_1 < \dots < x_{n-1} < x_n = b$, falls gilt $\{x_0, x_1, \dots, x_n\} \subseteq \{y_0, y_1, \dots, y_n\}$

$f: [a,b] \rightarrow \mathbb{R}$ ist \definition{Treppenfunktion}, falls $\exists$ Zerlegung $a=x_0 < x_1 < \dots < x_n = b$ von $[a,b]$ sodass $\forall k \in \{1, \dots, n\}$ die Restriktionen von $f$ auf $(x_{k-1}, x_k)$, $f|_{(x_{k-1},x_k)}$ konstant sind.

$\implies$ $f$ ist Treppenfunktion bzgl. der Zerlegung $a=x_0 < x_1 < \dots < x_n = b$.

Das \definition{Integral der Treppenfunktion} $f$ über $[a,b]$ ist die reelle Zahl $\int_{a}^{b}f(x)\, dx := \sum_{k=1}^{n}c_k(x_k - x_{k-1})$, wobei $c_k$ der Funktionswert von $f$ auf dem Intervall $(x_{k-1}, x_k)$ ist.

Sei $f:[a,b] \rightarrow \mathbb{R}$ eine Funktion. 
Dann ist die Menge der \definition{Obersummen $\mathcal{U}(f)$}$:=\{\int_{a}^{b}s(x)\, dx | s\in \mathcal{TF}; s \geq f\}$, 
respektive die Menge der \definition{Untersummen $\mathcal{L}(f)$}$:=\{\int_{a}^{b}s(x)\, dx | s\in \mathcal{TF}; s \leq f\}$
wobei $\mathcal{TF}$ die Menge der Treppenfunktionen bezeichnet.

\includegraphics[width=0.26\columnwidth]{04_funktionen/obersummen_untersummen}

Eine beschränkte Funktion $f:[a,b] \rightarrow \mathbb{R}$ nennt man \definition{Riemann-integrierbar}, falls gilt $\sup \mathcal{L}(f)=\inf \mathcal{U}(f)$.

In diesem Fall wird der Wert das \definition{Riemann-Integral} von $f$ $\int_{a}^{b} f \, dx := \sup \mathcal{L}(f)=\inf \mathcal{U}(f)$ genannt.
$a$ nennt man \definition{untere Integrationsgrenze}, $b$ \definition{obere Integrationsgrenze} und $f$ den \definition{Integranden.}

\theorem{Monotone Funktionen sind Riemann-integrierbar.}
\textit{Monotone Funktionen sind in Intervall $[a,b]$ beschränkt.}

\textit{Gilt auch, wenn Funktion auf manchen Teilintervallen monoton wachsend und auf manchen Teilintervallen monoton fallend ist.}

\theorem{Stetige Funktionen sind Riemann-integrierbar.}
\textit{müssen aber nicht differenzierbar sein (z.B. Knickstellen ok)!}

\textit{Gilt auch für Funktionen die \greenBox{in endlich vielen Stellen nicht monoton oder stetig} sind (wegen Aufteilung des Integrationsbereichs).}


\subsubsection{Rechenregeln}
Seien $f,g:[a,b] \rightarrow \mathbb{R}$ zwei (Riemann-)integrierbare Funktionen.
\begin{enumerate}
    \item \textbf{Linearität:} Seien zudem $\alpha, \beta \in \mathbb{R}$. Dann ist auch $\alpha f + \beta g$ integrierbar und es gilt
        $$\int_{a}^{b} (\alpha f + \beta g)(x) \,dx=\alpha \int_{a}^{b} f(x) \, dx + \beta \int_{a}^{b} g(x) \, dx$$
    \item \textbf{Monotonie:} Falls gilt $f \leq g$, so gilt auch
        $$\int_{a}^{b} f(x) \, dx \leq \int_{a}^{b} g(x) \, dx$$
    \item \textbf{Dreicksungleichung:} 
        $$|\int_{a}^{b} f(x) \, dx| \leq \int_{a}^{b} |f(x)| \, dx$$
        \textit{Vorstellung: Positionsänderung ist $\leq$ zurückgelegte Strecke}
    \item \textbf{Umkehrung der Integrationsrichtung:}
        $$\int_{a}^{b} f(x) \, dx = - \int_{b}^{a} f(x) \, dx$$
        \textit{Vorstellung: Wir laufen wieder zurück}
    \item \textbf{Aufteilung des Integrationsbereichs:}
        $$\int_{a}^{b} f(x) \, dx = \int_{a}^{c} f(x) \, dx + \int_{c}^{b} f(x) \, dx, a \leq c \leq b$$
        \textit{Insb. interessant, wenn Integrationsbereich symmetrisch an Nullpunkt. Z.B. bei ungeraden Funktionen im Intervall $[-1, 1]$ ist Integral $0$, bei geraden Funktionen 2x der Wert von $[0,1]$.}
\end{enumerate}
\textit{\textbf{Beweise für Rechenregeln:} "Was passiert für Treppenfunktionen?", dann Approximations-Idee, allgemeiner zeigen.}

\theorem{Hauptsatz der Integral- und Differentialrechung:}

Sei $f:[a,b]\rightarrow \mathbb{R}$ stetig. Dann ist für alle $C \in \mathbb{R}$ die Funktion $$F(x):=\int_{a}^{x}f(y)\,dy + C$$ eine Stammfunktion von $f$. Jede Stammfunktion int von dieser Form.
\textit{(Konstruktion von Stammfunktionen)}

\textit{
    \textbf{Beweis:} 
    $f$ stetig $\implies$ $f$ auf jedem Intervall $[a,x], x \in [a,b]$ integrierbar.
    Dann zeigen, dass $F$ Stammfunktion von $f$ ist (mit $\epsilon / \delta$-Stetigkeitsdefinition), $\lim_{x \rightarrow x_0} |\frac{F(x)-F(x_0)}{x-x_0} - f(x_0)|=\lim |F'(x_0)-f(x_0)|=0$ (einzeln zeigen für $x \in (x_0, x_0+ \delta) \cap [a,b]$ und $x \in (x_0-\delta, x_0) \cap [a,b]$).
}

Sei $F:[a,b]\rightarrow \mathbb{R}$ stetig differenzierbar. Dann gilt für alle $x \in [a,b]$
$$F(x)= \underbrace{F(a)}_{\textit{\small{"Startwert"}}} + \underbrace{\int_{a}^{x} F'(t) \, dt}_{\textit{\small{"Aufsummierung aller Veränderungseffekte"}}}$$

\textit{\textbf{Herleitung:} Gem. Hauptsatz $F(x)=\int_{a}^{x}F'(y)\, dy + C, C \in \mathbb{R}$.
Insbesondere: $F(a)=\underbrace{\int_{a}^{a}F'(y)\,dy}_{=0}+C$}

Sei $f:[a,b] \rightarrow \mathbb{R}$ stetig und $F:[a,b] \rightarrow \mathbb{R}$ Stammfunktion von $f$. Dann gilt
$$\int_{a}^{b}f(t)\,dt = F(x)|_a^b = F(b)-F(a)$$
\textit{(Berechnung von bestimmten Integralen via Stammfunktion)}

\greenBox{Verallgemeinerung für unstetige $f$}
Für eine Funktion $\tilde{f}(y)$ mit endlich vielen Unstetigkeitsstellen $\{x_1, \dots, x_n\}$ ist $F(x):=\int_{a}^{x}f(y)\,dy+C$ noch eine Stammfunktion für $x \in [a,b] \ \{x_1, \dots, x_n\}$

\greenBox{Eindeutigkeit der Stammfunktion}
Der Hauptsatz liefert alle möglichen Stammfunktion.

\textit{\textbf{Beweis:} Sei $G(x)$ eine weitere Stammfunktion, dann $H(x):=G(x)-\int_{a}^{x}f(y) \, dy$. Dann gilt $H'(x)=f(x) - f(x)=0 \forall x \in [a,b]$. Dann muss aber $H(x)$ eine Konstante sein, d.h. $H(x)=C, C \in \mathbb{R}$, dies bedeutet $C = G(x)- \int_{a}^{x}f(y)\,dy \iff G(x)= \int_{a}^{x}f(y)\,dy + C, C \in \mathbb{R}$}

\theorem{Partielle Integration}
Seien $f,g:[a,b] \rightarrow \mathbb{R}$ stetig differenzierbare Funktionen.
$$\int_{a}^{b} f(x) g'(x) \,dx = f(x)g(x)|_a^b-\int_{a}^{b}f'(x)g(x)\,dx$$
\textit{Folgt aus Hauptsatz und Produkteregel: $\int_{a}^{b}\frac{d}{dx} f(x)\cdot g(x)\,dx = \int_{a}^{b}(f'(x)\cdot g(x) + f(x) \cdot g'(x))\,dx = \int_{a}^{b}f'(x)\cdot g(x) \,dx + \int_{a}^{b}f(x) \cdot g'(x)\,dx=f \cdot g|^b_a=f(b)g(b)-f(a)g(a) \implies \int_{a}^{b}f(x)g'(x)\,dx = f(x) \cdot g(x) |^b_a - \int_{a}^{b}f'(x)\cdot g(x) \, dx$}

\theorem{Substitution, Version 1}
Seien $I,J \subset \mathbb{R}$ zwei Intervalle, $f : I \rightarrow J$ stetig differenzierbar und $g:J \rightarrow \mathbb{R}$ stetig. Dann gilt für alle $[a,b] \subset I$
$$\int_{a}^{b}g(f(x))f'(x)\,dx = \int_{f(a)}^{f(b)}g(y)\,dy$$

\textit{Weil wir nicht mehr nach $x$ integrieren sondern nach $y$ müssen wir auch die Integrationsgrenzen anpassen und "schauen, wie wir die $x$ in $y$ übersetzen können".}

\textit{Ggf. muss es etwas angepasst werden wenn bspw. eine Faktor als Abweichung bei der Ableitung besteht.}

{\color{gray}
\theorem{Substitution, Version 2}
Seien $I,J \subset \mathbb{R}$ Intervalle, $f:I\rightarrow J$ stetig differenzierbar und $g:J \rightarrow \mathbb{R}$ stetig. 
Für $[a,b]\subset I$ gelte $f'(x)\neq 0$ für alle $x \in [a,b]$ und sei $f^{-1}:[f(a),f(b)]\rightarrow \mathbb{R}$ Inverse von $f|_{[a,b]}$. 
Dann
$$\int_{a}^{b} g(f(x))\,dx = \int_{f(a)}^{f(b)}g(y)(f^{-1})'(y)\,dy$$
}

\subsection{Funktionenfolgen}

Es gibt auch Folgen $(f_n)_{n \in \mathbb{N}_0}$ von Funktionen $f_n: D \subset \mathbb{R} \rightarrow \mathbb{R}$.

\begin{exercise}
    Beispiel einer Funktionenfolge: $f_n(x)=x^n$, dann sind die Glieder $f_1(x)=x, f_2(x)=x^2, f_3(x)=x^3, ...$
\end{exercise}

\definition{Punktweise Konvergenz}
Sei $f : D \subset \mathbb{R} \rightarrow \mathbb{R}$ eine weitere Funktion. Dann konvergiert die Folge $(f_n)_{n \in \mathbb{N}_0}$ punktweise gegen $f$, falls für jedes $x \in D$ die reelle Folge $(f_n(x))_{n \in \mathbb{N}_0}$ gegen $f(x)$ konvergiert.

Wenn für jedes feste $x$ gilt $\lim_{n \rightarrow \infty} f_n(x)=f(x)$

$f$ heisst auch \definition{Punktweiser Grenzwert} der Folge $(f_n)_{n \in \mathbb{N}_0}$

\begin{exercise}
    Punktweise Konvergenz zeigen:

    \textit{Betrachte fixierten Punkt $x$, berechne Grenzwert $f(x)=lim_{n \rightarrow \infty} f_n(x)$. Wenn dieser Grenzwert für jedes $x$ im Definitionsbereich existiert, dann konvergiert $f_n$ punktweise gegen $f$.}
\end{exercise}

\definition{Gleichmässige Konvergenz}
Sei $f : D \subset \mathbb{R} \rightarrow \mathbb{R}$ eine weitere Funktion. Dann konvergiert die Folge $(f_n)_{n \in \mathbb{N}_0}$ gleichmässig gegen $f$, falls für jedes $\epsilon > 0$ ein Index $N$ existiert, sodasss für alle $n \geq N$ und für alle $x \in D$ gilt $|f_n(x) - f(x)|< \epsilon$.

\begin{exercise}
    Gleichmässige Konvergenz zeigen:

    \textit{Zuerst punktweise Konvergenz zeigen. Dann: Wir betrachten den grössten Fehler im Definitionsbereich $\sup_{x \in D} |f_n(x)-f(x)|$. Falls gilt $\lim_{n \rightarrow \infty} \sup_{x \in D} |f_n(x)-f(x)|=0$, dann ist Konvergenz gleichmässig.}
\end{exercise}

Eine \greenBox{Potenzreihe} konvergiert auf jedem kompakten Teilintervall des Konvergenzintervalls \greenBox{gleichmässig}.

\theorem{Mittelwertsatz der Integralrechnung}
% TODO!
Der Mittelwertsatz der Integralrechnung besagt, dass eine stetige Funktion auf einem geschlossenen Intervall mindestens einmal ihren durchschnittlichen (mittleren) Funktionswert annimmt. 
Geometrisch bedeutet dies, dass die Fläche unter der Kurve genau dem Flächeninhalt eines Rechtecks mit derselben Intervallbreite und der Höhe dieses Mittelwerts entspricht.